\begin{figure}[!ht]
\centering
% BASE PROMPT BOX (new; comes first)
\begin{tcolorbox}[
    colback=promptbg!100!white,
    sharp corners=south, 
    boxrule=0.25mm, 
    fonttitle=\bfseries, 
    title={\textbf{Seed Prompt for Second-Hop of Multi-Hop QA System}}, 
    width=\textwidth, 
    enhanced,
    drop shadow
    ]
\scriptsize
\setlength{\parskip}{0pt}
\setlength{\itemsep}{0pt}

Given the fields \texttt{question}, \texttt{summary\_1}, produce the fields \texttt{query}.
\end{tcolorbox}
% \vspace{0.5mm}

% ANNOTATED SECOND-HOP QUERY PROMPT BOX (original; comes second)
\begin{tcolorbox}[
    colback=promptbg!100!white, 
    sharp corners=south, 
    boxrule=0.3mm, 
    fonttitle=\bfseries, 
    title={\textbf{GEPA's Optimized Prompt for Second-Hop of Multi-Hop QA System, GPT-4.1 Mini}}, 
    width=\textwidth, 
    enhanced,
    drop shadow
    ]

% LEFT PANEL (ANNOTATED EXCERPT; STRICT SUBSET, HIGHLIGHTED STRATEGY CALL-OUTS)

\scriptsize
\textbf{You will be given two input fields: \texttt{question} and \texttt{summary\_1}.} \textbf{Your task:} Generate a new search query (\texttt{query}) \emph{optimized for the second hop} of a multi-hop retrieval system.

\begin{itemize}[left=0em]
    \item The original user question is typically complex and requires information from multiple documents to answer.
    \item The first hop query is the original question (used to retrieve initial documents).
    \item Your goal: generate a query to retrieve documents \emph{not} found in first hop but necessary to answer the question completely.
\end{itemize}

\vspace{1mm}
\textbf{Input Understanding:} \texttt{question} is the original multi-hop question posed by the user. \texttt{summary\_1} is a concise summary of information from a document retrieved in the first hop, which partially addresses the question.

\vspace{1mm}
\textbf{Purpose and Context:}
\begin{itemize}[left=0em]
    \item Your generated \texttt{query} aims to find the \emph{missing pieces} of information needed to fully answer the question. \ldots
    \item The query must retrieve relevant documents \emph{NOT} found in first hop \ldots for final answer extraction.
\end{itemize}

\vspace{1mm}
\textbf{Key Observations and Lessons:}
\begin{itemize}[left=0em]
    \item First-hop documents often cover one entity or aspect.
    \item Remaining relevant documents often involve connected or higher-level concepts mentioned in \texttt{summary\_1} but not explicitly asked in the original question. The \texttt{query} should target these \emph{missing}, but logically linked, documents.
    \item Avoid merely paraphrasing the original question or restating known facts from \texttt{summary\_1}.
    \item Infer what broader or related entities/concepts might provide the crucial missing information.
    \item For example:
    \begin{itemize}[left=1em]
        \item If \texttt{summary\_1} describes a population for a small civil parish, but the question wants the total population of the wider region, your query should target that wider region (e.g., ``Madeira archipelago population in 2011'').
        \item If \texttt{summary\_1} covers a song and the question asks for the album, target album-level documents.
    \end{itemize}
\end{itemize}

\vspace{1mm}
\textbf{How to Build the Query:}
\begin{itemize}[left=0em]
    \item Identify entities or topics mentioned in \texttt{summary\_1} that are related but different from first-hop documents.
    \item Reframe the query to explicitly mention these broader or related entities \emph{connected to the original question}.
    \item Include relevant key context from the question to maintain specificity, but shift focus to the missing piece.
    \item The goal is to retrieve documents that link or complement what was retrieved initially.
\end{itemize}

\vspace{1mm}
\textbf{Practical Strategy:}
\begin{itemize}[left=0em]
    \item Read the \texttt{summary\_1} carefully to spot references to bigger contexts or other entities not covered in the first hop.
    \item Ask: ``What entity or aspect does this summary hint at that could answer the original question but was not found yet?''
    \item Formulate a precise, focused factual query targeting that entity or concept to retrieve the missing documents.
\end{itemize}

\vspace{1mm}
\textbf{Output:}
\begin{itemize}[left=0em]
    \item Produce \texttt{query} as a clear, concise question or keyword phrase designed for efficient retrieval of second-hop documents.
    \item Ensure the query relates logically to the original question while targeting the broader or complementary knowledge identified in \texttt{summary\_1}. \ldots~Do not include the original question or simply rephrase it. Do not duplicate information already well-covered by the first hop
retrieval \ldots
\end{itemize}
\end{tcolorbox}

\caption{This figure shows an example prompt generated by GEPA for the second-hop document retrieval to be performed in a multi-hop question-answer system, along with the seed prompt it started with. Appendix~\ref{appendix:all_prompts} compares GEPA's prompts for all tasks with prompts generated by MIPROv2.}\label{fig:main_fig_prompt_example}
\end{figure}

% \begin{figure}[ht]
% \centering
% % BASE PROMPT BOX (new; comes first)
% \begin{tcolorbox}[
%     colback=promptbg!100!white,
%     sharp corners=south, 
%     boxrule=0.25mm, 
%     fonttitle=\bfseries, 
%     title={\textbf{Seed Prompt for Second-Hop of Multi-Hop QA System}}, 
%     width=\textwidth, 
%     enhanced,
%     drop shadow
%     ]
% \small
% Given the fields \texttt{question}, \texttt{summary\_1}, produce the fields \texttt{query}.
% \end{tcolorbox}
% % \vspace{0.5mm}

% % ANNOTATED SECOND-HOP QUERY PROMPT BOX (original; comes second)
% \begin{tcolorbox}[
%     colback=promptbg!100!white, 
%     sharp corners=south, 
%     boxrule=0.3mm, 
%     fonttitle=\bfseries, 
%     title={\textbf{GEPA's Optimized Prompt for Second-Hop of Multi-Hop QA System}}, 
%     width=\textwidth, 
%     enhanced,
%     drop shadow
%     ]

% % LEFT PANEL (ANNOTATED EXCERPT; STRICT SUBSET, HIGHLIGHTED STRATEGY CALL-OUTS)

% \small
% \textbf{You will be given two input fields: \texttt{question} and \texttt{summary\_1}.} \textbf{Your task:} Generate a new search query (\texttt{query}) \emph{optimized for the second hop} of a multi-hop retrieval system.

% \begin{itemize}[left=0em]
%     \item The original user question is typically complex and requires information from multiple documents to answer.
%     \item The first hop query is the original question (used to retrieve initial documents).
%     \item Your goal: generate a query to retrieve documents \emph{not} found in first hop but necessary to answer the question completely.
% \end{itemize}

% \vspace{1mm}
% \textbf{Input Understanding:} \texttt{question} is the original multi-hop question posed by the user. \texttt{summary\_1} is a concise summary of information from a document retrieved in the first hop, which partially addresses the question.

% \vspace{1mm}
% \textbf{Purpose and Context:}
% \begin{itemize}[left=0em]
%     \item Your generated \texttt{query} aims to find the \emph{missing pieces} of information needed to fully answer the question. \ldots
%     \item The query must retrieve relevant documents \emph{NOT} found in first hop \ldots for final answer extraction.
% \end{itemize}

% \vspace{1mm}
% \textbf{Key Observations and Lessons:}
% \begin{itemize}[left=0em]
%     \item First-hop documents often cover one entity or aspect.
%     \item Remaining relevant documents often involve connected or higher-level concepts mentioned in \texttt{summary\_1} but not explicitly asked in the original question. The \texttt{query} should target these \emph{missing}, but logically linked, documents.
%     \item Avoid merely paraphrasing the original question or restating known facts from \texttt{summary\_1}.
%     \item Infer what broader or related entities/concepts might provide the crucial missing information.
%     \item For example:
%     \begin{itemize}[left=1em]
%         \item If \texttt{summary\_1} describes a population for a small civil parish, but the question wants the total population of the wider region, your query should target that wider region (e.g., ``Madeira archipelago population in 2011'').
%         \item If \texttt{summary\_1} covers a song and the question asks for the album, target album-level documents.
%     \end{itemize}
% \end{itemize}

% \vspace{1mm}
% \textbf{How to Build the Query:}
% \begin{itemize}[left=0em]
%     \item Identify entities or topics mentioned in \texttt{summary\_1} that are related but different from first-hop documents.
%     \item Reframe the query to explicitly mention these broader or related entities \emph{connected to the original question}.
%     \item Include relevant key context from the question to maintain specificity, but shift focus to the missing piece.
%     \item The goal is to retrieve documents that link or complement what was retrieved initially.
% \end{itemize}

% \vspace{1mm}
% \textbf{Practical Strategy:}
% \begin{itemize}[left=0em]
%     \item Read the \texttt{summary\_1} carefully to spot references to bigger contexts or other entities not covered in the first hop.
%     \item Ask: ``What entity or aspect does this summary hint at that could answer the original question but was not found yet?''
%     \item Formulate a precise, focused factual query targeting that entity or concept to retrieve the missing documents.
% \end{itemize}

% \vspace{1mm}
% \textbf{Output:}
% \begin{itemize}[left=0em]
%     \item Produce \texttt{query} as a clear, concise question or keyword phrase designed for efficient retrieval of second-hop documents.
%     \item Ensure the query relates logically to the original question while targeting the broader or complementary knowledge identified in \texttt{summary\_1}. \ldots~Do not include the original question or simply rephrase it. Do not duplicate information already well-covered by the first hop
% retrieval \ldots
% \end{itemize}
% \end{tcolorbox}

% \caption{This figure shows an example prompt generated by GEPA for the second-hop document retrieval to be performed in a multi-hop question-answer system, along with the seed prompt it started with. Appendix~\ref{appendix:all_prompts} compares GEPA's prompts for all tasks with prompts generated by MIPROv2.}\label{fig:main_fig_prompt_example}
% \end{figure}

% \begin{figure}[ht]
% \centering
% % Given the fields `question`, `summary_1`, produce the fields `query`.
% \begin{tcolorbox}[
%     colback=promptbg!100!white, 
%     sharp corners=south, 
%     boxrule=0.3mm, 
%     fonttitle=\bfseries, 
%     title={\textbf{Annotated Second-Hop Query Generation Prompt (Multi-Hop Retrieval)}}, 
%     width=\textwidth, 
%     enhanced,
%     drop shadow
%     ]

% % LEFT PANEL (ANNOTATED EXCERPT; STRICT SUBSET, HIGHLIGHTED STRATEGY CALL-OUTS)

% \small
% \textbf{You will be given two input fields: \texttt{question} and \texttt{summary\_1}.} \textbf{Your task:} Generate a new search query (\texttt{query}) \emph{optimized for the second hop} of a multi-hop retrieval system.

% \begin{itemize}[left=0em]
%     \item The original user question is typically complex and requires information from multiple documents to answer.
%     \item The first hop query is the original question (used to retrieve initial documents).
%     \item Your goal: generate a query to retrieve documents \emph{not} found in first hop but necessary to answer the question completely.
% \end{itemize}

% \vspace{1mm}
% \textbf{Input Understanding:} \texttt{question} is the original multi-hop question posed by the user. \texttt{summary\_1} is a concise summary of information from a document retrieved in the first hop, which partially addresses the question.

% \vspace{1mm}
% \textbf{Purpose and Context:}
% \begin{itemize}[left=0em]
%     \item Your generated \texttt{query} aims to find the \emph{missing pieces} of information needed to fully answer the question. \ldots
%     \item The query must help retrieve other relevant documents \emph{NOT} found in the first hop that hold complementary or broader context necessary for final answer extraction.
% \end{itemize}

% \vspace{1mm}
% \textbf{Key Observations and Lessons:}
% \begin{itemize}[left=0em]
%     \item First-hop documents often cover one entity or aspect.
%     \item Remaining relevant documents often involve connected or higher-level concepts mentioned in \texttt{summary\_1} but not explicitly asked in the original question. The \texttt{query} should target these \emph{missing}, but logically linked, documents.
%     \item Avoid merely paraphrasing the original question or restating known facts from \texttt{summary\_1}.
%     \item Infer what broader or related entities/concepts might provide the crucial missing information.
%     \item For example:
%     \begin{itemize}[left=1em]
%         \item If \texttt{summary\_1} describes a population for a small civil parish, but the question wants the total population of the wider region, your query should target that wider region (e.g., ``Madeira archipelago population in 2011'').
%         \item If \texttt{summary\_1} covers a song and the question asks for the album, target album-level documents.
%     \end{itemize}
% \end{itemize}

% \vspace{1mm}
% \textbf{How to Build the Query:}
% \begin{itemize}[left=0em]
%     \item Identify entities or topics mentioned in \texttt{summary\_1} that are related but different from first-hop documents.
%     \item Reframe the query to explicitly mention these broader or related entities \emph{connected to the original question}.
%     \item Include relevant key context from the question to maintain specificity, but shift focus to the missing piece.
%     \item The goal is to retrieve documents that link or complement what was retrieved initially.
% \end{itemize}

% \vspace{1mm}
% \textbf{Practical Strategy:}
% \begin{itemize}[left=0em]
%     \item Read the \texttt{summary\_1} carefully to spot references to bigger contexts or other entities not covered in the first hop.
%     \item Ask: ``What entity or aspect does this summary hint at that could answer the original question but was not found yet?''
%     \item Formulate a precise, focused factual query targeting that entity or concept to retrieve the missing documents.
% \end{itemize}

% \vspace{1mm}
% \textbf{Output:}
% \begin{itemize}[left=0em]
%     \item Produce \texttt{query} as a clear, concise question or keyword phrase designed for efficient retrieval of second-hop documents.
%     \item Ensure the query relates logically to the original question while targeting the broader or complementary knowledge identified in \texttt{summary\_1}. \ldots~Do not include the original question or simply rephrase it. Do not duplicate information already well-covered by the first hop
% retrieval \ldots
% \end{itemize}
% \end{tcolorbox}
% \caption{Excerpt of \textbf{annotated second-hop retrieval query prompt}.}\label{fig:main_fig_prompt_example}
% \label{fig:annotated_secondhop_prompt}
% \end{figure}



% \begin{figure}[ht]
% \centering
% \begin{tcolorbox}[
%     colback=promptbg!100!white, 
%     sharp corners=south, 
%     boxrule=0.3mm, 
%     fonttitle=\bfseries, 
%     title={\textbf{Annotated Second-Hop Query Generation Prompt (Multi-Hop Retrieval)}}, 
%     width=\textwidth, 
%     enhanced,
%     drop shadow
%     ]

% % LEFT PANEL (ANNOTATED EXCERPT; STRICT SUBSET, HIGHLIGHTED STRATEGY CALL-OUTS)

% \small
% \textbf{You will be given two input fields: \colorbox{highlight}{\strut \texttt{question}} and \colorbox{highlight}{\strut \texttt{summary\_1}}.}

% \vspace{1mm}
% \textbf{Your task:} \colorbox{highlight}{\strut generate a new search query (\texttt{query}) \emph{optimized for the second hop}} of a multi-hop retrieval system.

% \begin{itemize}[left=0em]
%     \item The original user question is typically complex and requires information from multiple documents to answer.
%     \item The \colorbox{highlight}{\strut first hop query is the original question} (used to retrieve initial documents).
%     \item \colorbox{highlight}{\strut Your goal: generate a query} to retrieve \colorbox{highlight}{\strut additional relevant documents \emph{not} found in the first hop} but necessary to answer the question completely.
% \end{itemize}

% \vspace{1mm}
% \textbf{Input Understanding:}
% \begin{itemize}[left=0em]
%     \item \texttt{question} is the \colorbox{highlight}{\strut original multi-hop question posed by the user}.
%     \item \texttt{summary\_1} is a \colorbox{highlight}{\strut concise summary of information from a document retrieved in the first hop}, which partially addresses the question.
% \end{itemize}

% \vspace{1mm}
% \textbf{Purpose and Context:}
% \begin{itemize}[left=0em]
%     \item Your \colorbox{highlight}{\strut generated \texttt{query} aims to find the \emph{missing pieces} of information} needed to fully answer the question.
%     \item \ldots
%     \item \colorbox{highlight}{\strut The query must help retrieve other relevant documents \emph{NOT} found in the first hop} that hold complementary or broader context necessary for final answer extraction.
% \end{itemize}

% \vspace{1mm}
% \textbf{Key Observations and Lessons:}
% \begin{itemize}[left=0em]
%     \item First-hop documents \colorbox{highlight}{\strut often cover one entity or aspect}.
%     \item \colorbox{highlight}{\strut Remaining relevant documents often involve connected or higher-level concepts mentioned in \texttt{summary\_1} but not explicitly asked in the original question.}
%     \item The \colorbox{highlight}{\strut \texttt{query} should target these \emph{missing}, but logically linked, documents.}
%     \item \colorbox{highlight}{\strut Avoid merely paraphrasing the original question or restating known facts from \texttt{summary\_1}}.
%     \item \colorbox{highlight}{\strut Infer what broader or related entities/concepts might provide the crucial missing information.}
%     \item For example:
%     \begin{itemize}[left=1em]
%         \item If \texttt{summary\_1} describes a population for a small civil parish, but the question wants the total population of the wider region, your query should target that wider region (e.g., \colorbox{highlight}{\strut ``Madeira archipelago population in 2011''}).
%         \item If \texttt{summary\_1} covers a song and the question asks for the album, target album-level documents.
%     \end{itemize}
% \end{itemize}

% \vspace{1mm}
% \textbf{How to Build the Query:}
% \begin{itemize}[left=0em]
%     \item \colorbox{highlight}{\strut Identify entities or topics mentioned in \texttt{summary\_1} that are related but different from first-hop documents.}
%     \item \colorbox{highlight}{\strut Reframe the query to explicitly mention these broader or related entities \emph{connected to the original question}.}
%     \item Include relevant key context from the question to maintain specificity, but shift focus to the missing piece.
%     \item The goal is to \colorbox{highlight}{\strut retrieve documents that link or complement what was retrieved initially}.
% \end{itemize}

% \vspace{1mm}
% \textbf{Practical Strategy:}
% \begin{itemize}[left=0em]
%     \item \colorbox{highlight}{\strut Read the \texttt{summary\_1} carefully to spot references to bigger contexts or other entities not covered in the first hop}.
%     \item Ask: \colorbox{highlight}{\strut ``What entity or aspect does this summary hint at that could answer the original question but was not found yet?''}
%     \item \colorbox{highlight}{\strut Formulate a precise, focused factual query targeting that entity or concept to retrieve the missing documents}.
% \end{itemize}

% \vspace{1mm}
% \textbf{Output:}
% \begin{itemize}[left=0em]
%     \item Produce only the field \texttt{query} as a clear, concise question or keyword phrase designed for efficient retrieval of second-hop documents.
%     \item \colorbox{highlight}{\strut Ensure the query relates logically to the original question while targeting the broader or complementary knowledge identified in \texttt{summary\_1}.}
%     \item \colorbox{highlight}{\strut Do not include the original question or simply rephrase it}.
%     \item \colorbox{highlight}{\strut Do not duplicate information already well-covered by the first hop retrieval}.
% \end{itemize}

% \vspace{2mm}

% \textbf{Generalizable Lessons \& Strategies (for prompt designers):}
% \begin{itemize}[left=0em]
%     \item \colorbox{highlight}{\strut Multi-hop retrieval often requires \emph{bridging between retrieved factual content and broader context} not yet seen.}
%     \item \colorbox{highlight}{\strut Input summaries from earlier hops are crucial for clue extraction—teach annotators to look for entities, numbers, timeframes, or relationships that indicate scope gaps.}
%     \item \colorbox{highlight}{\strut Good second-hop prompt design should:}
%     \begin{itemize}[left=1em]
%         \item Encourage logical inference from partial answers rather than mechanical restatement.
%         \item Provide real-world or domain-specific examples (as above) to illustrate reframing strategies.
%         \item Clearly warn against uninformative/overlapping queries (“Do NOT…”).
%     \end{itemize}
% \end{itemize}

% \end{tcolorbox}
% \caption{Excerpt of \textbf{annotated second-hop retrieval query prompt}, with meta-level lessons for multi-hop QA prompt design. This summary is a strict subset of the full deployed task instruction.}
% \label{fig:annotated_secondhop_prompt}
% \end{figure}



%$%%%%%%%%%%%%%%%%%%%%%%%%%%%%%%%%%%%%%%%%%%%%%%$
% \begin{figure}[ht]
% \centering
% % \begin{tcolorbox}[
% %     colback=promptbg!100!white, 
% %     sharp corners=south, 
% %     boxrule=0.3mm, 
% %     fonttitle=\bfseries, 
% %     title={\textbf{Autonomously Generated Structured Summarization Prompt (GEPA)}}, 
% %     width=0.99\textwidth, 
% %     enhanced,
% %     drop shadow,
% %     sidebyside, 
% %     sidebyside align=top, 
% %     sidebyside gap=4mm,
% %     lefthand width=0.99\textwidth, 
% %     righthand width=0.01\textwidth
% %     ]
% \begin{tcolorbox}[
%     colback=promptbg!100!white, 
%     sharp corners=south, 
%     boxrule=0.3mm, 
%     fonttitle=\bfseries, 
%     title={\textbf{Example Prompt Generated by GEPA for HotpotQA Summarize Module}}, 
%     width=\textwidth, 
%     enhanced,
%     drop shadow
%     ]

% % LEFT PANEL (EXTRACTED PROMPT ITSELF, SHORTENED & HIGHLIGHTED)
% \small
% \textbf{You are the first-hop \colorbox{highlight}{\strut summarization module} in a multi-hop QA system. Your task is to generate a comprehensive, structured summary that:}

% \vspace{1mm}
% \begin{itemize}[left=0em]
% \item \textbf{Extracts direct answers} from the top retrieved passages.
% \item \textbf{Identifies missing or implied clues} that may require further retrieval
% \item \textbf{Synthesizes information}: combining explicit facts with domain knowledge or \textbf{logical inferences} to guide next steps.
% \end{itemize}

% \vspace{1mm}
% \textbf{Summary Structure}:
% \begin{itemize}[left=0em]
%     \item \textbf{Entity/Person Mention:} Clearly state the subject (\emph{full names, designations, ...}).
%     \item \textbf{Direct Answer:} Include \colorbox{highlight}{\strut explicit answers} from the passages (e.g., birth dates, team affiliations, or direct statements).
%     \item \textbf{Clues for Next Steps:} \colorbox{highlight}{\strut Signal missing information} (e.g., "Lance Rentzel's birth year is explicitly stated, but his exact birthplace is not; \textbf{need to search for 'Lance Rentzel birthplace'}").
%     \item \textbf{Domain-Specific Context:} Add relevant background (\emph{location, club history, etc.}).
% \end{itemize}

% \vspace{1mm}
% \textbf{Guidelines}:
% \begin{itemize}[left=0em]
%     \item \textbf{Do not omit any detail that could be relevant for follow-up queries} (e.g., team names, historical context).
%     \item \textbf{Prioritize clarity:} \colorbox{highlight}{\strut Separate direct answers from inferred clues}, e.g., using bullet points, subheadings\ldots
%     \item \textbf{Avoid assumptions not supported} by the passages;
    
%     \colorbox{highlight}{\strut if information is absent, explicitly state it is missing \& suggest precise search terms} (e.g., "Verify Wichita Dwight D. Eisenhower National Airport's tower status via FAA records").
%     \item \colorbox{highlight}{\strut Include quantifiable data:} (\emph{dates, numbers, etc.}) for precise comparisons.
%     \item \colorbox{highlight}{\strut Highlight connections} between entities (e.g., "Billy Truax \ldots were traded in 1970") to aid in cross-referencing.
% \end{itemize}

% \vspace{1mm}
% \textbf{Key Niche/Domain-Specific Insights}:
% \begin{itemize}[left=0em]\setlength\itemsep{1pt}
%     \item Birth dates are critical for age determination \ldots team affiliations (e.g., "traded in 1970") imply historical context.  
%     \item \ldots
% \end{itemize}

% \vspace{1mm}
% \textbf{Critical Additional Instructions}:
% \begin{itemize}[left=0em]\setlength\itemsep{1pt}
%     \item \textbf{Ensure All Retrieved Docs Are Represented}: Include all entities, and details from retrieved passages
%     \item \textbf{Signal Missing Links}: If connection b/w entities is implied but not explicitly stated suggest search terms to resolve it.
%     \item \textbf{Prioritize bridging concepts}: Highlight relation b/w entities (e.g., "Gary Pinkel coached Toledo in 1993 and holds the most wins in school history") to enable focused follow-up queries.
% \end{itemize}

% \vspace{1mm}
% \textbf{Example Format:}
% \vspace{-1mm}
% % \begin{lstlisting}[basicstyle=\ttfamily\tiny, backgroundcolor=\color{white}, frame=none]
% \begin{lstlisting}[basicstyle=\ttfamily, backgroundcolor=\color{white}, frame=none]
% Q: Which NFL player is younger, Billy Truax or Lance Rentzel?
% Entity/Person: Billy Truax (William Frederick Truax), Lance Rentzel (Thomas Lance Rentzel)
% Direct Answer:
% - Billy Truax: Born July 15, 1943.
% - Lance Rentzel: Born October 14, 1943.
% Clues for Next Steps: None; birth dates provided.
% Domain-Specific Context: Birth dates suffice to determine age difference.
% \end{lstlisting}


% \end{tcolorbox}
% \caption{Excerpt from prompt created by GEPA for summarize module to solve HotPotQA.}\label{fig:main_fig_prompt_example}
% \end{figure}
%%%%%%%%%%%%%%%%%%%%%%%%%%%%%%%%

% % RIGHT PANEL (ANNOTATIONS ABOUT WHY THIS PROMPT IS SO STRONG)

% \tcblower
% \footnotesize
% \begin{tcolorbox}[colback=annotbg, boxrule=0pt, sharp corners, left=0pt, right=0pt, top=0pt, bottom=0pt]
% \textbf{What Makes This Prompt Remarkable?}
% \begin{itemize}
%     \item \textbf{Explicit multi-level structure}: Separates direct, inferred, and missing content for robust downstream use.
%     \item \textbf{Automatic error signaling}: Instructs to explicitly state when needed information is missing and to propose actionable search terms to overcome gaps.
%     \item \textbf{Comprehensive detail preservation}: Ensures all entities, facts, and quantifiable data are surfaced without omission.
%     \item \textbf{Careful guidance for generalization:} Critical instructions prevent over-generalization and drive the system to stay evidence-based.
%     \item \textbf{Bridging concepts and next-step reasoning}: Promotes synthesis and relational reasoning for follow-up queries—enabling multi-hop QA.
%     \item \textbf{Full example outputs}: Concrete formats for output drive clarity and evaluation consistency.
% \end{itemize}
% \end{tcolorbox}


% \begin{instructbox}[\incode{HotpotQABench qwen3-8b summarize1.predict}]
% GEPA Prompt:
% \begin{lstlisting}[breakindent=0pt, breaklines=true, breakatwhitespace=true, frame=single, backgroundcolor=\color{blue!10}]
% You are the first-hop **summarization module** in a multi-hop QA system. Your task is to generate a **comprehensive, structured summary** that:  

% 1. **Extracts direct answers** from the top retrieved passages to address the question.  
% 2. **Identifies and highlights missing or implied clues** that may require further retrieval (e.g., entities, connections, or contextual details).  
% 3. **Synthesizes information** by combining explicit facts from the passages with domain-specific knowledge or logical inferences to guide subsequent steps.  

% ### **Summary Structure**  
% - **Entity/Person Mention**: Clearly state the subject (e.g., "Billy Truax", "Eintracht Braunschweig") and include **full names, titles, or official designations** (e.g., "Thomas Lance Rentzel", "Braunschweiger Turn- und Sportverein Eintracht von 1895 e.V.").  
% - **Direct Answer**: Include **explicit answers** from the passages (e.g., birth dates, team affiliations, or direct statements).  
% - **Clues for Next Steps**: Signal **missing information** (e.g., "Lance Rentzel's birth year is explicitly stated, but his exact birthplace is not; need to search for 'Lance Rentzel birthplace'").  
% - **Domain-Specific Context**: Add **relevant background** (e.g., "Eintracht Braunschweig is a German football club based in Braunschweig, Lower Saxony" or "NFL players' birth dates are critical for age comparisons").  

% ### **Guidelines**  
% - **Do not omit** any entity or detail from the retrieved passages that could be relevant for follow-up queries (e.g., team names, locations, or historical context).  
% - **Prioritize clarity** by **separating direct answers from inferred clues** (e.g., using bullet points, subheadings, or bolded labels).  
% - **Avoid assumptions** not supported by the passages; if information is absent, **explicitly state that it is missing** and suggest **precise search terms** (e.g., "Verify Wichita Dwight D. Eisenhower National Airport's tower status via FAA records").  
% - **Include quantifiable data** (e.g., "few thousand Stabyhouns exist globally", "born July 15, 1943") to enable precise comparisons.  
% - **Highlight connections** between entities (e.g., "Billy Truax and Lance Rentzel were traded in 1970") to aid in cross-referencing.  

% ### **Key Niche/Domain-Specific Insights**  
% - **NFL Player Comparison**: Birth dates are critical for age determination, and team affiliations (e.g., "traded in 1970") may imply historical context.  
% - **Airport Classification**: "Non-towered" status is explicitly stated in some passages (e.g., "non-towered public airport"), while others require inference (e.g., "major commercial airports typically have towers").  
% - **Football Club Context**: Clubs like Eintracht Braunschweig require background on their location, league, and history (e.g., "based in Braunschweig, Lower Saxony").  
% - **Quantifiable Data**: Use exact dates, numbers, or rankings (e.g., "few thousand Stabyhouns exist globally") to enable precise comparisons.  

% ### **Critical Additional Instructions**  
% - **Ensure All Retrieved Documents Are Represented**: Explicitly include all entities, titles, and details from the retrieved passages (e.g., full names, film titles, and specific roles).  
% - **Signal Missing Links**: If a connection between entities is implied but not explicitly stated (e.g., "Nancy Steiner worked on *The Lovely Bones*"), flag this as a potential gap and suggest search terms to resolve it.  
% - **Prioritize Bridging Concepts**: Highlight relationships between entities (e.g., "Gary Pinkel coached Toledo in 1993 and holds the most wins in school history") to enable focused follow-up queries.  
% - **Avoid Overgeneralization**: Only include domain-specific context that is either explicitly stated in the passages or directly inferable (e.g., "major commercial airports typically have towers" is acceptable, but "airports with fewer than 10,000 passengers are non-towered" is not unless stated).  

% ### **Example Format**  
% For the question *"Which NFL player is younger, Billy Truax or Lance Rentzel?"*:  
% - **Entity/Person Mention**: Billy Truax (William Frederick Truax), Lance Rentzel (Thomas Lance Rentzel)  
% - **Direct Answer**:  
%   - **Billy Truax**: Born July 15, 1943.  
%   - **Lance Rentzel**: Born October 14, 1943.  
% - **Clues for Next Steps**: None required; birth dates are explicitly provided.  
% - **Domain-Specific Context**: Birth dates are sufficient to determine age difference within the same year.  

% For the question *"Which is a non-towered airport, Wichita Dwight D. Eisenhower National Airport or Montrose Regional Airport?"*:  
% - **Entity/Person Mention**: Wichita Dwight D. Eisenhower National Airport, Montrose Regional Airport  
% - **Direct Answer**:  
%   - **Montrose Regional Airport**: "non-towered public airport" (passage 3).  
%   - **Wichita Dwight D. Eisenhower National Airport**: No explicit mention of tower status; inferred as **towered** (typical for major commercial airports).  
% - **Clues for Next Steps**: Verify Wichita's tower status via FAA records or additional sources (e.g., "Wichita Dwight D. Eisenhower National Airport tower status").  
% - **Domain-Specific Context**: Non-towered airports lack a control tower, relying on pilot communication (passage 4). Major commercial airports like Wichita usually have towers.  

% **Tip:** When summarizing, don’t just compress; synthesize—include both direct answers and clues required for the system’s next steps. Always explicitly state if a retrieved document’s content is missing critical information, and provide actionable search terms to address gaps.
% \end{lstlisting}
% \end{instructbox}